\noindent
Real-time analysis of high-velocity event streams is a cornerstone requirement for modern monitoring, telemetry, and analytics platforms, where systems frequently process millions of records per second. Extracting relevant aggregated signals such as distinct element counts, exact item frequencies, and heavy hitters in an exact manner requires storing state proportional to the traffic volume. This causes $O(N)$ memory bloat and unacceptable latency delays for real-time alerting. Consequently, exact methods are often infeasible in production systems with strict memory budgets and Service Level Agreements (SLAs).

This project presents a highly scalable, adaptive system architecture that relies on probabilistic data structures to provide bounded-error estimations. We combine three primary probabilistic sketches --- Count-Min Sketch (CMS) for point-query frequency estimates, HyperLogLog (HLL) for cardinality approximations, and the Sliding-Window Misra--Gries (MG) algorithm for tracking trending heavy hitters under concept drift over continuous sliding windows. 

In the first phase of this development, we effectively established the core stream-processing foundation. A robust stream ingestion pipeline was erected using Apache Kafka to simulate Zipfian-distributed heavy-tailed web traffic. Independent, highly-scalable worker nodes process this stream and independently compute multi-sketch estimates. To rigorously assess performance ceilings, both standard Python asyncio workers and hardware-accelerated C++ workers (utilizing \textit{std::countl\_zero} and XXH64 hashing) were constructed and integrated. The distributed sketch state is synchronized safely using an in-memory Redis cluster that serves as a high-speed federated state store. Finally, we deployed a responsive FastAPI-based REST API backing an interactive web dashboard to monitor system health, explore live sketch estimations, and track P99 latencies in real time. 

The successful culmination of this static-parameter phase demonstrates high-throughput real-time sketching capability and lays out the infrastructure necessary. In the subsequent phase, this baseline will be enhanced through the integration of an Adaptive Runtime Controller and a Load-Aware Query Router that dynamically scales precision parameters and query distribution dynamically, directly reacting to changing node memory pressure and latency SLA constraints. 
